\documentclass[10pt]{article}
\usepackage[margin=1in]{geometry}
\usepackage[T1]{fontenc}
\usepackage{lmodern,amsmath,amssymb,amsthm,microtype,enumitem}
\usepackage[colorlinks=true,allcolors=blue]{hyperref}
\newtheorem{theorem}{Theorem}
\newcommand{\Pp}{\mathbb P}
\newcommand{\E}{\mathbb E}
\title{A designed tree model with a sharp congestion--fragmentation transition}
\author{Alec Kriebel\\\small Independent researcher\\
\small\href{https://orcid.org/0009-0001-9320-500X}{ORCID: 0009-0001-9320-500X}}
\date{21 October 2026\quad (version 1.0; unrefereed preprint)}
\begin{document}
\maketitle
\begin{abstract}
We construct a bounded-degree physical tree model in which the normalized
marginal admitted demand drops from one to zero while the spare-capacity graph
changes from the entire tree to isolated vertices. Capacities arise from the
unique cheapest fixed-topology enlargement supporting doubled initial traffic.
Each of the \(k^2\) cross-root streams has an independent, fixed-law exponential
growth rate. An all-cuts estimate confines both changes, with probability
tending to one, to an interval of width \(16\sqrt{\log k/k}\) around mean
doubled demand. Both outer-phase conclusions hold for every optimizer.
This is a specialized existence response to Aldous's conceptual congestion
toy-model request. The width is an upper bound; no critical profile or universal
traffic law is asserted.
\end{abstract}

\section{The request and the model}
Aldous's problem page \cite{A} asks for a tractable model connecting a rapid
decline in the marginal amount of extra traffic admitted with fragmentation
of the graph of edges having spare capacity. One suggested design mechanism
minimizes the cost of increasing capacities to support doubled current demand.
The accompanying index classifies the request as conceptual, rather than as
one uniquely formulated conjecture. We implement that mechanism in an explicit
family. Its topology, traffic pattern and growth process are part of the model.

Let \(T_H\) be the complete rooted binary tree of height \(H\ge2\), with
\(k=2^{H-1}\) leaves in each root subtree and \(N=4k-1\) vertices.
Write \(L,R\) for its left and right leaf sets. There is one commodity
from each \(i\in L\) to each \(j\in R\), initially of size one; all other
demands vanish. Each commodity uses its unique simple physical path.
For an edge \(e\), let \(k_e\) count its descendant leaves.
Its initial load and capacity are \(kk_e\). Give every capacity increment
a strictly positive linear price and require all doubled demands to fit.
The unique cheapest resulting capacities are
\begin{equation}\label{capacity}
 C_e=2kk_e.
\end{equation}
Indeed, the paths force these lower bounds separately, and the bounds are
simultaneously feasible. Positive prices force the componentwise minimum.
Terminal capacities grow as \(2k\); the two root-edge capacities are
\(2k^2\). Each finite network is then held fixed.

Sample independent \(X_{ij}\sim\operatorname{Exp}(1)\), and for \(t\ge1\) put
\[
 d_{ij}(t)=1+(t-1)X_{ij},\qquad S_k=\sum_{i,j}X_{ij}>0\quad\hbox{a.s.}
\]
These are quenched random slopes, not fresh samples at successive times.
Each stream has mean \(t\) and variance \((t-1)^2\), independently of \(k\).
Thus this is a specialization different from the linear-variance example
on the source page.

Let \(F_k(t)\) be the maximum of \(\sum_{ij}y_{ij}\), with
\(0\le y_{ij}\le d_{ij}(t)\), subject to physical capacities. Define
\[
 r_k(t)=F'_{k,+}(t)/S_k.
\]
The denominator is the realized offered-demand slope, not initial demand
\(k^2\). For each optimizing admitted matrix let \(E_k(t)\) contain exactly
the edges with loads strictly below capacity, on the full vertex set of
\(T_H\); write \(M_k(t)\) for its largest component size.
Inside the transition this graph may depend on the optimizer. The two
phase conclusions below quantify over every optimizer. Component-size
convergence applies to any measurable selection of an optimizer;
the outer-phase event itself is uniform over all optimizers.

\section{The joint transition}
\begin{theorem}\label{main}
Put
\[
 \delta_k=8\sqrt{\frac{\log k}{k}},\qquad
 \epsilon_k=2k^{-31}+\frac{2k^{-2}}{(1-k^{-2})^2}
                   +\left(\frac4{k^2}\right)^k.
\]
For all binary \(k\) large enough that \(\delta_k<1\), with probability at
least \(1-\epsilon_k\), both statements hold simultaneously:
\begin{enumerate}[label=(\roman*),leftmargin=2em]
\item For every \(1\le t\le2-\delta_k\), all demand is admitted,
\(r_k(t)=1\), and every optimizer has \(E_k(t)=T_H\), hence \(M_k(t)=N\).
\item For every \(t\ge2+\delta_k\), \(F_k(t)=2k^2\), \(r_k(t)=0\),
and every optimizer has edgeless \(E_k(t)\), hence \(M_k(t)=1\).
\end{enumerate}
Consequently, for fixed \(1\le t<2\), \(r_k(t)\to1\) and \(M_k(t)/N\to1\)
in probability; for fixed \(t>2\) both tend to zero.
For every fixed \(t\ge1\), including \(t=2\),
\[
 F_k(t)/k^2\longrightarrow\min(t,2)\quad\hbox{in probability}.
\]
No limit of \(r_k(2)\), critical component-size profile, or optimal
transition-width order is asserted.
\end{theorem}

\subsection{Exact deterministic reduction}
A left leaf's load is its row sum \(\sum_jy_{ij}\); a right leaf's load
is its column sum. Each other load is the sum of the corresponding
descendant terminal loads. Thus \eqref{capacity} is equivalent exactly to
\begin{equation}\label{constraints}
 \sum_j y_{ij}\le2k,\qquad \sum_i y_{ij}\le2k,\qquad
 0\le y_{ij}\le d_{ij}(t).
\end{equation}
Internal constraints introduce no further restriction.

The source--\(L\)--\(R\)--sink network has capacities \(2k\),
\(d_{ij}(t)\), \(2k\) on its three arc layers. Max-flow/min-cut
\cite{FF} gives the finite formula
\begin{equation}\label{cut}
 F_k(t)=\min_{\substack{A\subseteq L\\B\subseteq R}}
 \left[2k(k-|A|+|B|)+|A|(k-|B|)
             +(t-1)\sum_{i\in A,j\notin B}X_{ij}\right].
\end{equation}
It is a finite minimum of nondecreasing affine functions, hence is
continuous, concave and piecewise affine. Its right derivative is the
smallest active slope. With \(\mathcal I_k(t)\) the active cuts, including
all ties,
\begin{equation}\label{marginal}
 r_k(t)=\min_{(A,B)\in\mathcal I_k(t)}
             \frac{\sum_{i\in A,j\notin B}X_{ij}}{S_k}\in[0,1].
\end{equation}
These formulas do not assert efficient exhaustive enumeration or an
elementary finite-\(k\) expectation.

If all offered row and column sums are strictly less than \(2k\), full
admission is feasible and is the unique admitted matrix attaining the offered
total: every entry is individually upper bounded. Every physical edge
then has strict slack. Conversely, \(F_k(t)\le2k^2\) always. Equality forces
every row and column of every optimizer to equal \(2k\), saturating every
physical edge. These observations establish the optimizer quantifiers.

\subsection{Exponential tails and the lower phase}
For a sum \(Z\) of \(n\) independent \(\operatorname{Exp}(1)\) variables,
\(\E e^{\lambda Z}=(1-\lambda)^{-n}\) for \(\lambda<1\).
Exponential Markov with \(\lambda=1-1/a\) gives
\[
 \Pp(Z\ge na)\le e^{-nI(a)}\ (a>1),\qquad
 \Pp(Z\le na)\le e^{-nI(a)}\ (0<a<1),\quad I(a)=a-1-\log a.
\]
For \(0<\delta<1\),
\begin{equation}\label{rate}
 I((1-\delta)^{-1})\ge\delta^2/2,\qquad
 I((1+\delta)^{-1})\ge\delta^2/8.
\end{equation}
Both sides vanish at zero. The left-hand derivatives in \(\delta\) are
\(\delta/(1-\delta)^2\ge\delta\) and
\(\delta/(1+\delta)^2\ge\delta/4\), proving \eqref{rate} by integration.

At \(t=2-\delta\), an offered row sum is
\(k+(1-\delta)\sum_jX_{ij}\). It reaches \(2k\) only if that exponential
sum reaches \(k/(1-\delta)\); the same holds for a column.
A union bound over \(2k\) events, without joint independence, gives
\[
 \Pp(\hbox{some offered row or column reaches }2k)
       \le2k e^{-k\delta^2/2}.
\]
At \(\delta=\delta_k\) this is \(2k^{-31}\). Otherwise all rows and columns
have strict slack at the endpoint and at all earlier \(t\ge1\).
The endpoint slack also persists in a right neighborhood. The deterministic
observations prove (i), including its endpoint right derivative.

\subsection{All upper-phase cuts}
At \(t=2+\delta\), all cuts in \eqref{cut} are at least \(2k^2\)
exactly when
\[
 \sum_{i\in A,j\notin B}d_{ij}(t)\ge2k(|A|-|B|).
\]
Cuts with \(|A|\le|B|\) are automatic. For the others put
\(D=R\setminus B\), \(a=|A|\), \(d=|D|\), so \(a+d>k\).
Since
\[
 ad-k(a+d-k)=(k-a)(k-d)\ge0,
\]
failure implies
\[
 ad+(1+\delta)\sum_{A\times D}X_{ij}<2k(a+d-k)\le2ad,
 \qquad \sum_{A\times D}X_{ij}<ad/(1+\delta).
\]
Its probability is at most \(e^{-cad}\), where
\(c=I((1+\delta)^{-1})\ge\delta^2/8\). Different cuts need not be independent.

Put \(q=\min(a,d)\), \(\ell=k-\max(a,d)\). Then \(0\le\ell<q\) and
\(ad=q(k-\ell)\). If \(q\le\lfloor k/2\rfloor\), the number of choices
for fixed \(q,\ell\), allowing interchange, is at most
\(2\binom{k}{q}\binom{k}{\ell}\). Summing over \(\ell<q\) bounds it by
\(2q k^{2q}\), whereas \(ad\ge qk/2\). For \(q>k/2\), \(ad>k^2/4\)
and there are at most \(4^k\) choices in total. The probability that any
cut fails is therefore at most
\[
 \sum_{q=1}^{\lfloor k/2\rfloor}2q k^{2q}e^{-cqk/2}
                     +4^k e^{-ck^2/4}.
\]
At \(\delta_k\), \(c\ge8\log(k)/k\), bounding this by
\[
 \sum_{q\ge1}2q k^{-2q}+\left(\frac4{k^2}\right)^k
   =\frac{2k^{-2}}{(1-k^{-2})^2}+\left(\frac4{k^2}\right)^k.
\]
When all cuts pass, \(F_k(2+\delta_k)=2k^2\). Increasing demands retain a
maximizing admitted flow for all later times. The upper bound \(2k^2\)
then makes throughput constant, its right derivative zero, and every
edge saturated for every optimizer. Combining the two failure bounds
proves (i)--(ii).

\subsection{Limits}
We have \(\E(S_k/k^2)=1\) and
\(\operatorname{Var}(S_k/k^2)=k^{-2}\).
The fixed-\(t\) conclusions away from two follow from (i)--(ii),
\(\epsilon_k\to0\), and \(\delta_k\to0\). At two, on the same good event,
\[
 k^2+(1-\delta_k)S_k=F_k(2-\delta_k)
                    \le F_k(2)\le2k^2.
\]
This gives the remaining throughput limit without exchanging limits and
derivatives. The failure bounds are summable along binary \(k\); hence the
outer-phase events also hold eventually almost surely under any common
coupling, by the first Borel--Cantelli lemma.

\section{Context, scope and reproducibility}
Capacitated random transportation and flow-cut analysis are classical.
Hassin--Zemel \cite{HZ} and Karp--Motwani--Nisan \cite{KMN,KMNreport} are
essential antecedents, alongside max-flow/min-cut and elementary concentration.
The highly optimized tolerance paradigm is prior work \cite{CD}, as are
congestion-induced fragmentation and shortest-path percolation \cite{KR}.
The latter sequentially consumes edges for individual requests, rather than
optimizing simultaneous admitted demand under fixed capacities.
The contribution here is the explicit cheapest-design tree construction and
its simultaneous normalized-marginal and every-optimizer spare-connectivity
theorem, not a new general transport-feasibility principle.

The sufficient exponent separation in \cite[Theorem 1]{HZ} does not cover
the present balanced near-capacity regime. The directed theorem in the actual
1988 report \cite[Section 4]{KMNreport} fixes bounded integer capacities
and a positive constant mean surplus; Section 5 separately proposes a
vanishing-surplus extension. That report has not been equated with an inspected
1993 final edition: the final journal body \cite{KMN} was unavailable.
These comparisons do not exclude every possible derivation from prior work.

The physical graph is sparse but traffic is dense, capacities increase with
size, and internal constraints reduce to terminal constraints. Only the
fixed-topology design alternative is implemented. No adaptive capacity growth,
generic cyclic-network theorem, linear-variance process, critical profile or
universal critical law is supplied. Throughput saturates; the marginal drops.
The width is a sufficient bound without a matching lower bound. Aldous's
2005 discussion \cite[slide 5]{Atalk} already contrasts easy but unrealistic
tree analysis with preferred mean-field models. Our tree has no rerouting
and does not remove that explanatory limitation.
The suitability of one example as the uniquely \emph{right} conceptual model
is not a mathematical conclusion.

A dated bounded audit inspected 16 primary source bodies and recorded
52 literature searches, locating no earlier exact construction and theorem
in that corpus. The final-edition gap above, unpublished work and incomplete
indexing remain limitations; no first-discovery or global novelty certificate
is asserted. The source and theorem comparisons are in the supplement.

The proof is analytic; finite controls do not establish rare-event asymptotics.
The original 24-file two-turn packet is preserved unchanged. The supplement's
new replay derivative verifies those frozen files and requires exact active
second-turn output bytes. It admits bounded absolute drift only at 40 specified
older floating diagnostics, together with explicit runtime and public-input
projection rules; all other structure and controls remain exact.
Python 3.9.6, NumPy 2.0.2 and SciPy 1.13.1 were tested.
The strict historical wrapper fails at roundoff in that environment and is
not advertised as portable. Exact policies, guard tests and limitations are
provided; other platform/library combinations are not guaranteed.

\paragraph{AI involvement and review status.}
AI tools were extensively involved in mathematical research, model and proof
development, computational verification, literature search, manuscript
preparation and adversarial auditing. This is an unrefereed preprint.
Automated independent-agent reviews are not independent human peer review;
the manuscript has not undergone formal independent human peer review.

\begin{thebibliography}{9}
\bibitem{A}
D. Aldous, \emph{Congestion in networks: SOC, HOT and the
percolation-fragmentation phase transition}, problem page posted May 2003,
update marked February 2018, accessed 21 October 2026.
\href{https://www.stat.berkeley.edu/~aldous/Research/OP/congestion.html}{Primary problem};
\href{https://www.stat.berkeley.edu/~aldous/Research/OP/index.html}{conceptual problem index}.
\bibitem{FF}
L. R. Ford, Jr. and D. R. Fulkerson, Maximal flow through a network,
\emph{Canadian Journal of Mathematics} \textbf{8} (1956), 399--404.
\href{https://doi.org/10.4153/CJM-1956-045-5}{doi:10.4153/CJM-1956-045-5}.
\bibitem{HZ}
R. Hassin and E. Zemel, Probabilistic analysis of the capacitated
transportation problem, \emph{Mathematics of Operations Research}
\textbf{13}(1) (1988), 80--89.
\href{https://doi.org/10.1287/moor.13.1.80}{doi:10.1287/moor.13.1.80}.
\bibitem{KMN}
R. M. Karp, R. Motwani and N. Nisan, Probabilistic analysis of network
flow algorithms, \emph{Mathematics of Operations Research}
\textbf{18}(1) (1993), 71--97.
\href{https://doi.org/10.1287/moor.18.1.71}{doi:10.1287/moor.18.1.71}.
\bibitem{KMNreport}
R. M. Karp, R. Motwani and N. Nisan, \emph{Probabilistic analysis of network
flow algorithms}, Berkeley technical report CSD-88-392, December 1988.
\href{https://www2.eecs.berkeley.edu/Pubs/TechRpts/1988/CSD-88-392.pdf}{Actual report}.
\bibitem{CD}
J. M. Carlson and J. Doyle, Complexity and robustness,
\emph{Proceedings of the National Academy of Sciences}
\textbf{99}(Suppl. 1) (2002), 2538--2545.
\href{https://doi.org/10.1073/pnas.012582499}{doi:10.1073/pnas.012582499}.
\bibitem{KR}
M. Kim and F. Radicchi, Shortest-path percolation on random networks,
\emph{Physical Review Letters} \textbf{133} (2024), 047402.
\href{https://doi.org/10.1103/PhysRevLett.133.047402}{doi:10.1103/PhysRevLett.133.047402}.
\bibitem{Atalk}
D. Aldous, \emph{Flows through random networks}, Ottawa talk, June 2005.
\href{https://www.stat.berkeley.edu/~aldous/Talks/Ottawa.pdf}{Primary slides}.
\end{thebibliography}
\end{document}
